%% kapa: exam=1 round=36 year=2025 period=2 form=A subject=accounting subject_ko=회계학 seq=36 number=76
%% source: Q-Net 감정평가사 기출문제 (cst00309) artlSeq=5231525 file=raw/1008/5231525/2025년 제36회 감정평가사 1차시험 2교시 문제지 원본.pdf
%% conversion: kordoc 4.24.51 -> tools/gen1.py (auto); review=proofread-sonnet; flags=none
\begin{question}{76}
전부원가계산, 변동원가계산 및 초변동원가계산에 관한 설명으로 옳지 않은 것은?
\begin{choices}
  \item 초변동원가계산은 직접재료원가만을 제품원가에 포함하고 나머지 제조원가는 모두 기간비용으로 처리한다.
  \item 변동원가계산은 이익계획 및 의사결정 목적에 유용하도록 원가를 변동원가와 고정원가로 분류하고 공헌이익을 보고한다.
  \item 전부원가계산하의 영업이익은 판매량뿐만 아니라 생산량의 변화에도 영향을 받지만, 변동원가계산하의 영업이익은 판매량에 의해서만 영향을 받는다.
  \item 전부원가계산과 변동원가계산은 수익과 비용의 대응원칙에 부합되는 원가계산방법으로 외부보고 및 조세 목적을 위해서 일반적으로 인정되는 방법이다.
  \item 초변동원가계산은 판매량이 일정한 경우에 생산량이 증가할수록 기간비용화되는 변동 가공원가가 증가하여 영업이익이 감소되므로 불필요한 재고의 누적을 방지하는 효과가 변동원가계산보다 크다.
\end{choices}
\end{question}
\answer{4}
